% !TEX root = ../main.tex
\chapter{MLX Runtime: Graph-Eager Hybrid}
\label{ch:mlx_runtime}

\Lucid's GPU backend is not implemented against \Metal\ directly. It
is implemented against \MLX, Apple's open-source machine-learning
framework for Apple silicon, which Apple's machine-learning research
group released in late 2023 \cite{hannun_mlx_2023}. The choice to sit
on top of \MLX\ rather than to wrap \Metal\ directly is the most
consequential single architectural decision in \Lucid: it determines
the shape of the dispatch layer, the cost model the autograd engine
must respect, the format of the intermediate representation in
\code{lucid.compile}, and even the contents of the wheel that
\code{pip install lucid} delivers (which carries an \MLX\ dependency
constraint of \(\geq 0.31\)). This chapter explains what \MLX\ is,
how its execution model works, and how \Lucid\ uses it.

The exposition is organised around four observations. First,
\MLX\ is a \emph{lazy graph} runtime: ordinary operations build an
unevaluated computation graph, and a separate \code{eval} call
materialises results. This is unusual among contemporary ML
frameworks and shapes nearly every consequence that follows. Second,
\MLX's primitive set is intentionally closer to NumPy's than to
\refframework's; some operations \Lucid\ exposes are constructed in
the Python composite layer (\chref{ch:composite_ops}) rather than
delegated to \MLX. Third, \MLX's runtime exposes a small set of
\emph{streams} that abstract over CPU and GPU execution; \Lucid's
device model uses these streams as its bedrock. Fourth, \MLX\ has
quirks --- some inherited from NumPy, some specific to Apple silicon
--- that the framework above must guard against; we catalogue the
two that have most affected \Lucid's correctness.

\section{Overview}
\label{sec:mlx_runtime:overview}

\MLX\ is a research-driven framework with a small core team and a
fast release cadence; the API surface that \Lucid\ targets has been
broadly stable since version \(0.10\), with consistent additions but
few removals. The framework has the following high-level shape:

\begin{itemize}
  \item A C++ core implementing the array primitive, the lazy graph,
    the function transformations, and the kernel set.
  \item A Python binding (\code{mlx.core}) that exposes the C++
    surface as Python objects.
  \item Companion Python modules (\code{mlx.nn}, \code{mlx.optimizers},
    \code{mlx.utils}) that build higher-level constructs on top of
    the core. \Lucid\ does not use these companion modules; we build
    our own equivalents.
  \item Metal kernels (in MSL) compiled at install time or on first
    use, plus a small CPU fallback path used when the array is on
    the CPU stream.
\end{itemize}

The piece \Lucid\ targets is the C++ core, accessed through the
Python binding for most user-facing operations and through
\code{lucid/\_C/backend/gpu/MlxBackend.\{h,cpp\}} for engine-side
operations. The Lucid C++ engine links against \code{libmlx.dylib}
at build time and dispatches GPU work through MLX's API.

\subsection{What MLX is, what MLX is not}
\label{ssec:mlx_runtime:what_is}

\MLX\ is not a deep-learning framework in the same sense that
\refframework\ is. It does not own a module class hierarchy, it does
not own an autograd engine with arbitrary backward functions, it
does not own an optimiser zoo. It is closer to a unified-memory
analogue of JAX's \code{jax.numpy} layer: an array library with
function transformations (\code{vmap}, \code{grad}, \code{jit}) that
the user composes manually.

This narrow scope is the property that makes \MLX\ a good substrate
for \Lucid. We do not duplicate \MLX's array implementation; we
delegate to it. We do not reimplement \MLX's kernels; we call them.
But we add (i) a parameter-and-module abstraction that \MLX\ lacks,
(ii) an autograd engine that handles \refframework-compatible custom
\code{Function} classes (\chref{ch:custom_functions}), (iii) a
dispatcher that selects between \MLX, \Accelerate, and \MPSGraph
depending on operation and device, (iv) an extensive model zoo, and
(v) a graph-capture compiler (\chref{ch:compile_design}). The
relationship is layered: \Lucid\ on \MLX, not \Lucid\ alongside
\MLX.

\subsection{Why MLX, not Metal directly}
\label{ssec:mlx_runtime:why_mlx}

A reasonable alternative would have been for \Lucid\ to wrap \Metal\
directly, in the way \refframework\ wraps CUDA directly. Three
considerations tipped the decision toward \MLX:

\begin{enumerate}
  \item \textbf{Kernel coverage.} \MLX\ ships a mature library of
    Metal kernels for elementwise, reduction, indexing, GEMM,
    convolution, FFT, and reduction-with-gather operations. Writing
    these from scratch would have been a multi-year effort and would
    have produced kernels that converge on \MLX's anyway.
  \item \textbf{Lazy graph as a free optimisation tier.} Because
    \MLX\ is lazy, sequences of operations fuse opportunistically.
    \Lucid\ benefits from these fusions without owning a separate
    optimiser. The trade-off is the synchronisation discipline of
    \secref{sec:mlx_runtime:execution}, which we accept.
  \item \textbf{Maintenance bandwidth.} A two-engineer team cannot
    maintain a full Metal kernel library and a full ML framework. By
    leaning on \MLX, \Lucid\ spends its engineering budget on the
    layers above the kernel (autograd, modules, training, model zoo)
    where the user-visible features are.
\end{enumerate}

The cost is a dependency we do not control: when \MLX\ ships a
breaking API change in a new minor release, \Lucid\ must adapt.
\secref{sec:mlx_runtime:versioning} discusses the discipline by
which we manage this.

\section{The Graph-Eager Hybrid Execution Model}
\label{sec:mlx_runtime:execution}

The single most important property of \MLX\ for the rest of this
book is its execution model. \MLX\ is \emph{lazy}: an ordinary
operation does not run when it is invoked. It is added to an internal
computation graph, and a result placeholder is returned. The
placeholder appears to the caller as an array, with shape and dtype
metadata available immediately, but with no backing buffer. The
backing buffer materialises only when \emph{evaluation} is forced.

\subsection{The lazy graph}
\label{ssec:mlx_runtime:lazy_graph}

\figref{fig:mlx_runtime:lazy_graph} sketches the model. Operations
construct nodes in a DAG; the DAG is materialised by an explicit
\code{eval} call (or by an implicit read of the bytes through, e.g.,
the bridge of \chref{ch:unified_memory}).

\begin{figure}[t]
\centering
\begin{adjustbox}{max width=\textwidth, center}
\begin{tikzpicture}[
  user/.style={draw, rounded corners=2pt, font=\small\ttfamily,
               fill=lucid.code.bg, minimum height=22pt,
               minimum width=2.4cm, inner sep=4pt},
  prim/.style={draw, rounded corners=2pt, font=\small\sffamily,
               fill=lucid.info.bg, minimum height=28pt,
               minimum width=3cm, inner sep=4pt, align=center},
  buf/.style={draw, rounded corners=2pt, font=\small\sffamily,
              fill=lucid.ok.bg, minimum height=42pt,
              minimum width=2.8cm, inner sep=4pt, align=center},
  arrow/.style={->, >=stealth, thick, color=lucid.accent.dark,
                shorten >=2pt, shorten <=2pt},
  ev/.style={->, >=stealth, very thick, color=lucid.warn.fg,
             dashed, shorten >=4pt, shorten <=4pt},
  header/.style={font=\small\bfseries, color=lucid.muted},
]
  % Column headers
  \node[header] at (-5.5, 4.2) {User code};
  \node[header] at (-0.5, 4.2) {Lazy graph (DAG)};
  \node[header] at ( 5.5, 4.2) {Backing buffer};

  % User-visible variables (left column)
  \node[user] (var_a) at (-5.5, 3.2) {a};
  \node[user] (var_b) at (-5.5, 2.2) {b};
  \node[user] (var_c) at (-5.5, 1.2) {c = a + b};
  \node[user] (var_d) at (-5.5, 0.2) {d = c * 2};

  % Graph nodes (middle column)
  \node[prim] (p_a) at (-0.5, 3.2) {leaf};
  \node[prim] (p_b) at (-0.5, 2.2) {leaf};
  \node[prim] (p_c) at (-0.5, 1.2) {primitive\_add};
  \node[prim] (p_d) at (-0.5, 0.2) {primitive\_mul};

  \draw[arrow] (p_a.east) -- ++(0.6,0) |- (p_c.east);
  \draw[arrow] (p_b.east) -- ++(0.4,0) |- (p_c.east);
  \draw[arrow] (p_c) -- (p_d);

  % Materialised buffer (right column, appears on eval)
  \node[buf] (buf_d) at (5.5, 0.2) {Buffer for d\\(materialised)};

  % Eval trigger
  \node[font=\small\ttfamily, color=lucid.warn.fg, align=center]
       (eval) at (0, -1.8)
    {mx.eval(d)\\or d.numpy()};
  \draw[ev] (eval.east) to[out=0, in=-90]
    node[font=\footnotesize\itshape, midway, below right, fill=white,
         inner sep=2pt] {walk DAG, schedule}
    (buf_d.south);
\end{tikzpicture}
\end{adjustbox}
\caption[MLX lazy graph execution.]{MLX's graph-eager hybrid
execution model. User code constructs arrays; each operation
appends a primitive node to the lazy DAG without computing
anything. The buffer materialises only when an explicit
\code{eval} or implicit read forces evaluation, at which point
MLX walks the DAG backward from the target, schedules primitives
on the appropriate stream, and writes the result into a backing
\code{MTL::Buffer}.}
\label{fig:mlx_runtime:lazy_graph}
\end{figure}

When the user writes
\begin{lstlisting}[style=lucid,language=LucidPython]
import mlx.core as mx
a = mx.array([1.0, 2.0, 3.0])
b = mx.array([4.0, 5.0, 6.0])
c = a + b           # no computation happens here
d = c * 2           # still no computation
\end{lstlisting}
the variable \code{d} holds an array object whose shape is
\((3,)\), whose dtype is \code{float32}, and whose underlying
buffer does not yet exist. The graph held by \MLX\ records:
\code{c = primitive\_add(a, b); d = primitive\_multiply(c, 2)}. Both
operations are \emph{deferred}.

The structure of this graph is a directed acyclic graph (DAG) of
\emph{primitives}, where each primitive is a leaf operation (an
elementwise add, a matmul, a reduction) and each node has zero or
more input arrays and one or more output arrays. An array is itself
a tuple of (primitive that produced it, position in that primitive's
output list, shape, dtype, strides).

\subsection{Evaluation triggers}
\label{ssec:mlx_runtime:eval_triggers}

Evaluation is forced by any of the following:

\begin{enumerate}
  \item An explicit \code{mx.eval(d)} or \code{d.eval()} call.
  \item A read of the array's bytes through \code{d.data<T>()}
    (which is what \code{numpy.asarray}, \code{tolist}, and \Lucid's
    \code{SharedStorage::from\_mlx} all do indirectly).
  \item A print or repr operation.
  \item An assignment into an \code{mx.array} via item assignment
    (\code{d[0] = 7.0}), which evaluates the array before
    in-place modification.
\end{enumerate}

Each of these triggers \MLX\ to walk the DAG backward from the
target, schedule primitives onto the appropriate stream, and block
until the result is materialised.

\subsection{Why lazy by default}
\label{ssec:mlx_runtime:lazy_why}

The lazy default is unusual; most ML frameworks (\refframework\
chief among them) are eager by default. The argument for laziness on
Apple silicon is concrete:

\begin{itemize}
  \item Apple silicon's GPU has high kernel-launch overhead relative
    to its kernel cost, especially for short kernels. Fusing adjacent
    kernels eliminates launches that would otherwise dominate the
    total time.
  \item The unified memory model makes intermediates cheap to fuse
    away. There is no allocation cost to be paid for fusion across
    primitives.
  \item Lazy evaluation enables operator-level optimisations
    (constant folding, broadcast collapsing, transpose absorption)
    without a separate compiler tier.
\end{itemize}

The cost is that the framework above \MLX\ must reason about when
evaluation happens. \Lucid\ has had two incidents in which a missed
evaluation produced a memory leak (\chref{ch:unified_memory}
Section~\ref{ssec:unified_memory:lazy_leak}) and one in which a
deferred evaluation produced a write-after-read race in an in-place
update. The fix in every case has been to insert a strategically
placed \code{eval} at a synchronisation point.

\subsection{Asynchronous evaluation}
\label{ssec:mlx_runtime:async_eval}

\MLX\ added \code{async\_eval} in version \(0.10\), which submits the
graph for evaluation without blocking. The caller may continue to
build the next portion of the graph while the GPU works on the
submitted portion; subsequent eager reads block as usual.

\Lucid\ uses \code{async\_eval} in two places. First, the optimiser
step (\chref{ch:optimizers}) evaluates the updated parameters
asynchronously so that the next forward pass can begin construction
in parallel. Second, the backward engine
(\chref{ch:autograd_engine}) submits accumulated gradients
asynchronously when their consumer is ready.

The benefit, on the training loop of a 100M-parameter Transformer,
is roughly \(5\)\,ms per backward pass on an M3 Max --- a small but
real fraction of total step time. The fix landed in 3.3.0 as part of
the AMP rollout; see \chref{ch:amp_integration}.

\section{The Array Primitive}
\label{sec:mlx_runtime:array}

The central data type in \MLX\ is \code{mlx::core::array} (Python:
\code{mlx.core.array}). It is the unit of every operation; it is
what \Lucid's \code{TensorImpl} wraps (\chref{ch:tensor_model}). We
describe its structure and the interface that the rest of \Lucid\
relies on.

\subsection{Structure}
\label{ssec:mlx_runtime:array_structure}

Conceptually, an \code{mlx::core::array} is:

\begin{lstlisting}[style=lucidcpp,language=C++]
struct array {
  Shape           shape;     // dimensions
  Strides         strides;   // element strides per dimension
  Dtype           dtype;     // element type
  Primitive*      primitive; // producer (null for leaf)
  std::vector<array> inputs; // upstream arrays
  size_t          out_index; // which output of primitive
  std::shared_ptr<Buffer> buffer; // materialised storage, or null
  // ... auxiliary fields for lazy evaluation
};
\end{lstlisting}

The fields that matter for \Lucid:

\begin{itemize}
  \item \code{shape}, \code{dtype} are always known and used by
    \Lucid's metadata layer to populate \code{Tensor.shape} and
    \code{Tensor.dtype} without forcing evaluation.
  \item \code{strides} encode lazy view operations
    (\code{transpose}, \code{slice}, non-trivial \code{reshape}) and
    must \emph{not} be passed naively to a CPU consumer, because
    \code{data<T>()} returns the bytes in their pre-view layout.
    This is the source of the bridge contiguity discipline of
    \chref{ch:unified_memory}.
  \item \code{buffer} is the materialised storage, null until
    evaluation. The pointer to its contents is what
    \code{data<T>()} returns.
  \item \code{primitive} and \code{inputs} define the DAG that
    \code{eval} walks.
\end{itemize}

\subsection{Dtype and conversion}
\label{ssec:mlx_runtime:array_dtype}

\MLX\ supports the following dtypes:

\begin{itemize}
  \item Floating point: \code{float16}, \code{float32}, \code{bfloat16}
    (since 0.5), \code{float64} (CPU stream only; not supported on
    GPU because \Metal\ has no native FP64).
  \item Integer: \code{int8}, \code{int16}, \code{int32},
    \code{int64}, and unsigned variants.
  \item \code{bool}, \code{complex64}.
\end{itemize}

\Lucid\ exposes the same set, modulo two carve-outs. First, \Lucid's
\code{float64} dtype is only valid on \code{device=`cpu'}; an attempt
to construct a \code{float64} tensor on \code{device=`metal'} raises
\code{LucidError(``float64 not supported on Metal'')}. Second,
\Lucid\ does not currently expose unsigned integer types in the
top-level surface, because \refframework\ parity favours signed
integers and the user demand for unsigned is small.

Conversion between dtypes uses \MLX's \code{astype} primitive, which
is lazy like any other. A common pitfall (one that has burned
\Lucid\ in 3.0.3) is that \code{astype} followed by a CPU bridge
must \code{eval} before the bridge, because \code{astype} is itself
deferred; the bridge's \code{contiguous().eval()} sequence does
this correctly only if it sees \code{astype} as the head of the
lazy chain.

\subsection{The data\textless T\textgreater\ accessor}
\label{ssec:mlx_runtime:array_data}

The most performance-critical method on \code{mlx::core::array} for
\Lucid\ is \code{data<T>()}. It returns a raw pointer to the
buffer's contents, typed as \code{T*}. The pointer is valid as long
as the array is live (and, transitively, as long as any \MLX\ array
that shares storage with it is live). The pointer is to bytes in
the array's stored layout, which is \emph{not} the same as the
strided layout if the array has been produced by a lazy view
operation.

The full discipline for using \code{data<T>()} from \Lucid\ is:

\begin{enumerate}
  \item Make the array contiguous via \code{mlx::core::contiguous(a)}.
  \item Evaluate via \code{eval(a)}.
  \item Access via \code{a.data<T>()}.
  \item Hold the array (a \code{shared\_ptr<array>}) for the
    lifetime of the CPU access.
\end{enumerate}

The \code{SharedStorage} class
(\chref{ch:unified_memory} Section~\ref{ssec:unified_memory:bridge_iface})
encapsulates this discipline. No \Lucid\ code outside that class
calls \code{data<T>()} directly.

\section{Streams and Devices}
\label{sec:mlx_runtime:streams}

\MLX's execution unit is a \emph{stream}: an ordered queue of
primitives that share a target device and execute in submission
order. There are two streams in the default configuration: the CPU
stream and the GPU stream. A primitive submitted to one stream
executes on that stream's device; primitives on different streams
execute concurrently and synchronise only when an array's bytes are
read or when an explicit barrier is inserted.

\subsection{The default device}
\label{ssec:mlx_runtime:default_device}

\MLX\ exposes a global default device through \code{mx.set\_default\_device()}.
The default is the GPU on machines that have one (i.e., all Apple
silicon), and the CPU on machines that do not. \Lucid\ does not use
the global default; the dispatcher (\chref{ch:backend_dispatch})
explicitly selects the stream per operation based on the tensor's
\code{device} attribute. This ensures that \Lucid's behaviour is
deterministic regardless of any global state \MLX\ may carry.

\subsection{Stream-aware primitives}
\label{ssec:mlx_runtime:stream_primitives}

Every \MLX\ primitive accepts an optional \code{stream=} or
\code{device=} argument that overrides the default for that
operation. \Lucid's dispatcher uses this to enforce per-op device
placement: a \code{cpu} tensor's operation is dispatched with
\code{stream=mx.cpu}, a \code{metal} tensor's operation with
\code{stream=mx.gpu}. The two streams may run concurrently, but
\Lucid's autograd engine enforces a happens-before relation at the
tensor boundary so that downstream operations see consistent state.

\subsection{The linalg carve-out}
\label{ssec:mlx_runtime:linalg_carveout}

There is one carve-out that is worth highlighting: linear-algebra
primitives (\code{cholesky}, \code{qr}, \code{svd}, \code{eigh},
\code{solve}, \code{lstsq}) are implemented in \MLX\ only on the CPU
stream. \Metal\ does not currently provide LAPACK-equivalent kernels;
\MLX\ uses \Accelerate's LAPACK under the hood, which is itself
CPU-only.

\Lucid's linalg API therefore has an unusual property: even when the
user has chosen \code{device=`metal'}, a \code{lucid.linalg.cholesky}
call dispatches on the CPU stream. The result is wrapped as a
metal-device tensor only after the linalg primitive returns. The
overhead is one \code{SharedStorage} bridge transfer per call, which
is negligible relative to the cost of the decomposition itself for
matrices larger than \(64 \times 64\). This is the \emph{linalg
carve-out} referenced throughout the book
(\chref{ch:linalg}, \chref{ch:backend_dispatch}); it is the only
operation in \Lucid\ for which device placement is partly fictive.

\section{The Primitive Set}
\label{sec:mlx_runtime:primitives}

\MLX's primitive set occupies a middle ground between NumPy and
\refframework. The naming convention favours NumPy
(\code{transpose}, not \code{permute}; \code{concatenate}, not
\code{cat}; \code{stack} but not \code{vstack}); the broadcasting
semantics are NumPy's; the reduction signatures
(\code{axis=}, \code{keepdims=}) are NumPy's.

\subsection{Elementwise and reduction}
\label{ssec:mlx_runtime:elementwise}

The elementwise set is comprehensive: arithmetic, comparison, logical
and bitwise operations, transcendental functions
(\code{exp}, \code{log}, \code{sin}, \code{cos}, \code{tanh},
\code{erf}, \code{erfinv}), rounding
(\code{floor}, \code{ceil}, \code{round}, \code{trunc}). \Lucid\
delegates nearly every elementwise primitive to \MLX\ directly,
exposing them with \refframework-style names through the dispatcher.

The reduction set covers \code{sum}, \code{mean}, \code{max},
\code{min}, \code{prod}, \code{var}, \code{std}, \code{any},
\code{all}, with per-axis and global variants. \code{argmax} and
\code{argmin} are also primitives. \code{logsumexp} is a primitive
since 0.18 (\Lucid\ implemented it as a composite before that).

\subsection{Shape manipulation}
\label{ssec:mlx_runtime:shape}

Reshape, transpose, broadcast-to, expand-dims, squeeze, swapaxes,
flatten, ravel, split, concatenate, stack. The interesting subset is
the lazy ones: \code{transpose} and \code{reshape} on most shapes
return a non-materialised array with adjusted strides (the
stride-trap discussed above). \code{concatenate} and \code{stack}
necessarily allocate.

\subsection{Indexing and gather/scatter}
\label{ssec:mlx_runtime:indexing}

\MLX\ supports advanced indexing through array indices (\code{a[idx]})
and through the \code{take} and \code{take\_along\_axis} primitives.
The scatter family is split into two primitives with similar names
but different semantics:

\begin{itemize}
  \item \code{scatter\_add(a, indices, updates)}: NumPy-style
    \code{np.add.at(a, indices, updates)} semantics. Multiple
    index lists, one per source dimension.
  \item \code{scatter\_add\_axis(a, axis, indices, updates)}:
    \refframework-style \code{scatter\_add\_(dim, indices, updates)}
    semantics. Single index list along one axis.
\end{itemize}

Passing one's indices to the wrong variant produces an immediate
shape-mismatch failure (\code{``Number of index arrays does not
match number of dimensions''}). \Lucid's gather/scatter wrappers
(\chref{ch:indexing}) defensively use the correct variant based on
the call shape; a vault note (\code{engine-mlx-scatter-axis-vs-multiaxis})
captures the diagnostic that surfaced this distinction the first
time.

\subsection{Linear algebra and FFT}
\label{ssec:mlx_runtime:linalg_fft}

The linalg primitives live in \code{mlx.core.linalg} and operate on
the CPU stream as discussed in
\secref{ssec:mlx_runtime:linalg_carveout}. The FFT primitives
(\code{fft}, \code{ifft}, \code{rfft}, \code{irfft}, with their
multi-dimensional variants) live in \code{mlx.core.fft} and run on
both streams; \Lucid's \code{lucid.fft} module
(\chref{ch:fft}) wraps them with consistent normalisation
conventions.

\subsection{Convolution and neural-network primitives}
\label{ssec:mlx_runtime:nn_prims}

\code{conv1d}, \code{conv2d}, \code{conv3d}, and their transposed
counterparts are primitives. Pooling, batch-norm-style fused kernels,
and attention's softmax fused with a scale are not first-class
primitives but rather compositions in the \code{mlx.nn} module that
\Lucid\ does not use; \Lucid's \code{nn.functional} layer
(\chref{ch:linear_conv_norm}) calls the core conv primitives
directly and composes the rest in its own kernels.

\section{Function Transformations}
\label{sec:mlx_runtime:transforms}

\MLX\ inherits from JAX a small set of function-level transformations
that operate on Python functions of arrays and return new functions
with new behaviours. The list is short: \code{grad}, \code{vmap},
\code{jit} (also called \code{mlx.compile}), \code{vjp},
\code{jvp}, and a few utilities.

\Lucid\ does not, for the most part, expose these transformations to
the user. The reason is parity: \refframework's user-facing API
expresses differentiation through the \code{Module.backward} and
\code{tensor.backward} pattern, not through a functional
\code{grad}. \Lucid\ therefore implements its own autograd engine
(\chref{ch:autograd_engine}) on top of \MLX's primitives, and
\MLX's \code{grad} sits unused.

The two transformations \Lucid\ does use are \code{vmap} and
\code{jit}, but only indirectly:

\begin{itemize}
  \item \code{vmap} is exposed in \Lucid\ as
    \code{lucid.func.vmap} (\chref{ch:higher_order}), where it
    composes onto \Lucid\ tensors. The implementation delegates to
    \MLX's \code{vmap} on the underlying arrays.
  \item \code{jit} (i.e., \code{mlx.compile}) is \emph{not} the
    mechanism behind \code{lucid.compile}. \Lucid's compile path
    (\chref{ch:compile_design}) emits onto \MPSGraph\ rather than
    \MLX's compile; the reasons are detailed in that chapter.
\end{itemize}

\section{Lucid's Use of MLX}
\label{sec:mlx_runtime:lucid_use}

We can now describe how \Lucid\ sits on top of \MLX. The integration
has three layers.

\subsection{The MlxBackend interface}
\label{ssec:mlx_runtime:mlx_backend}

\code{MlxBackend} (in \code{lucid/\_C/backend/gpu/}) is the C++ class
that mediates between \Lucid's dispatcher and \MLX. Its interface is
roughly:

\begin{lstlisting}[style=lucidcpp,language=C++]
class MlxBackend : public IBackend {
public:
  TensorImpl allocate(Shape, Dtype) override;
  TensorImpl unary(OpKind, const TensorImpl& x) override;
  TensorImpl binary(OpKind, const TensorImpl& a, const TensorImpl& b) override;
  TensorImpl reduce(OpKind, const TensorImpl& x, AxisSet, bool keepdims) override;
  TensorImpl matmul(const TensorImpl& a, const TensorImpl& b) override;
  // ... and so on for every primitive category
};
\end{lstlisting}

Each method translates a \Lucid\ operation request into the
corresponding \MLX\ primitive call, with the \code{stream} argument
fixed to \code{mx.gpu}. The return value is a new \code{TensorImpl}
that wraps the \MLX\ array returned by the primitive.

\subsection{Dispatch into MLX from the engine}
\label{ssec:mlx_runtime:dispatch}

The dispatcher (\chref{ch:backend_dispatch}) routes a per-op request
to one of the backends based on the tensor's device. For
\code{device=`metal'}, the route is \code{MlxBackend} unconditionally,
with one exception (the MPSGraph per-op carve-out introduced in 3.4;
\chref{ch:mpsgraph_dispatch}). For \code{device=`cpu'}, the route is
\code{AccelerateBackend} \emph{except} for the linalg carve-out,
where it is also \code{MlxBackend} but with \code{stream=mx.cpu}.

The Python-side dispatcher (\code{lucid/\_dispatch.py}) is symmetric
but operates on Python \code{Tensor} objects. It performs the
\code{\_wrap} / \code{\_unwrap} between \code{Tensor} and
\code{TensorImpl}, then defers to the C++ dispatcher for the actual
op routing.

\subsection{Bridge from MLX to Lucid Tensor}
\label{ssec:mlx_runtime:bridge}

The final layer is the bridge from \code{mlx::core::array} to
\Lucid's \code{TensorImpl}. The bridge wraps the array in a
\code{TensorImpl} that holds the array as a \code{shared\_ptr},
inherits the array's shape and dtype, and computes the tensor's
strides from the array's strides. Autograd metadata
(\code{requires\_grad}, \code{grad\_fn}) is allocated separately
because \MLX\ has no notion of these.

When the inverse direction is needed --- a \code{TensorImpl} that
came from elsewhere needs to be passed to an \MLX\ primitive ---
the bridge unwraps the \code{shared\_ptr<array>} that the
\code{TensorImpl} holds. This is always free; the array was
constructed on the \MLX\ side originally.

\section{Quirks and Pitfalls}
\label{sec:mlx_runtime:quirks}

Two \MLX\ behaviours have been the source of bugs in \Lucid's
history; both are documented in the engine vault, both have shaped
the framework's defensive coding.

\subsection{data\textless T\textgreater\ ignores strides}
\label{ssec:mlx_runtime:data_strides}

The first quirk has been discussed in detail in
\chref{ch:unified_memory}: \code{mlx::core::array::data<T>()} returns
a pointer to bytes in the array's stored layout, not the strided
layout. A lazy transpose of a \(4 \times 8\) array of dtype
\code{float32} returns an \(8 \times 4\) array whose
\code{data<T>()} pointer points to the original \(4 \times 8\)
contiguous block.

The discipline is enforced by \code{SharedStorage} and is
non-negotiable: \code{contiguous} then \code{eval} then
\code{data<T>()}. The vault note
\code{engine-mlx-data-ignores-strides} is the source-of-truth for
the regression that surfaced this in 3.0 series.

\subsection{Two scatter variants}
\label{ssec:mlx_runtime:two_scatters}

The second quirk is the existence of two primitives with similar
names and different semantics: \code{scatter\_add} (NumPy-style
multi-axis index arrays) and \code{scatter\_add\_axis}
(\refframework-style single-axis indices along one axis). A
\Lucid\ wrapper that picks the wrong variant produces a runtime
error of the form ``Number of index arrays does not match number of
dimensions of input''. The fix in 3.0 series was to standardise on
\code{scatter\_add\_axis} for the \refframework-compatible API and
to use \code{scatter\_add} only in NumPy-bridge code paths. The
vault note \code{engine-mlx-scatter-axis-vs-multiaxis} captures the
investigation.

\subsection{Axis vs.\ axes nomenclature}
\label{ssec:mlx_runtime:axis_axes}

\MLX\ uses \code{axis=} (singular) for some reductions and
\code{axes=} (plural) for others. The convention follows NumPy:
operations that accept exactly one axis use \code{axis=}; operations
that accept a tuple of axes (or a single axis as a tuple) use
\code{axes=}. \Lucid\ standardised on \code{dim=} and \code{dims=}
internally for \refframework\ parity, and the dispatcher rewrites
between conventions on each call. The Phase~7 engine rename retro
(\code{retro-axis-to-dim-rename}) records the design.

\section{Version Compatibility and the Wheel Constraint}
\label{sec:mlx_runtime:versioning}

\MLX\ is a young framework on a fast cadence. \Lucid\ pins a minimum
version (\code{mlx >= 0.31}) and tracks releases as they ship. The
constraint discipline has three parts.

\subsection{API stability per minor version}
\label{ssec:mlx_runtime:api_stability}

\MLX's minor releases (0.18, 0.20, etc.) sometimes introduce
breaking changes. The standard pattern is that \MLX\ deprecates an
API in version \(N\) and removes it in version \(N+2\). \Lucid's
build pins the minimum supported \MLX\ version conservatively and
the maximum loosely. The wheel that \code{pip install lucid} fetches
declares \code{mlx>=0.31,<1.0}, which gives \MLX\ room to ship
minor releases that \Lucid\ has not been tested against.

\subsection{Platform constraint: macOS \(\geq 26.0\)}
\label{ssec:mlx_runtime:macos_constraint}

The wheel constraint is tied to \MLX's own wheel constraint. As of
\MLX\ 0.31, the published wheel is tagged \code{macosx\_26\_0\_arm64},
which means it requires macOS 26.0 (Tahoe) and the arm64 architecture.
There is no Linux wheel; there is no Intel-Mac wheel; there is no
older-macOS wheel. \Lucid\ inherits these constraints and matches
them with its own build target (\code{MACOSX\_DEPLOYMENT\_TARGET=26.0}).

The \code{mlx-metal} split package (introduced in \MLX\ 0.20)
separates the Metal kernel binaries from the core Python package
to reduce wheel size. \Lucid's build links against the combined
package and does not need to know about the split.

\subsection{Per-version conditional code}
\label{ssec:mlx_runtime:per_version}

A small number of \Lucid\ code paths are guarded by \MLX\ version
checks. The pattern is:

\begin{lstlisting}[style=lucid,language=LucidPython]
from packaging.version import Version
import mlx.core as mx

if Version(mx.__version__) >= Version("0.18"):
    out = mx.logsumexp(x, axis=axis)
else:
    out = _composite_logsumexp(x, axis=axis)
\end{lstlisting}

These conditional paths are exceptional; the policy is to delete the
fallback once the minimum supported \MLX\ version exceeds the
introduction point. \Lucid's current minimum (\(0.31\)) elides
roughly half a dozen such fallbacks that were present in earlier
\Lucid\ releases.

\section{Summary and Forward References}
\label{sec:mlx_runtime:summary}

\MLX\ is the GPU-side substrate on which \Lucid\ is built. Its lazy
graph model gives \Lucid\ a free fusion tier and forces a
synchronisation discipline that the framework above must respect.
Its array primitive is the storage type that \Lucid's \code{Tensor}
wraps. Its stream model is the basis for \Lucid's device routing,
with the one carve-out that linear algebra is dispatched on the CPU
stream regardless of user device choice. Its quirks
(\code{data<T>()} stride-ignorance, two scatter variants, axis
nomenclature drift) have been the source of historical bugs that
have shaped \Lucid's defensive layers.

The next chapter, \chref{ch:accelerate}, takes the CPU-side
counterpart to \MLX --- the \Accelerate\ framework --- and describes
how its BLAS, LAPACK, vDSP, vForce, and BNNS sublayers compose into
the \code{AccelerateBackend} that \Lucid's dispatcher targets when
the user has chosen \code{device=`cpu'}.
\chref{ch:metal_mpsgraph} then closes \partref{part:hardware} by
treating \Metal\ and \MPSGraph\ directly, since both are needed for
the graph-capture compiler. After that, the book moves to the
software architecture proper.

\begin{lucidnote}
For day-to-day use, the relevant \MLX\ facts are: (1) \MLX\ is lazy,
so \code{eval} is the synchronisation primitive; (2) \MLX's array
is the storage layer \Lucid\ delegates to; (3) device placement in
\Lucid\ corresponds to stream selection in \MLX, with the linalg
carve-out. The remaining detail in this chapter is required when
debugging the bridge boundary or extending the dispatcher.
\end{lucidnote}
